\documentclass[11pt]{article}

\usepackage[margin=1in]{geometry}
\usepackage{amsmath,amssymb,amsthm,mathtools}
\usepackage[T1]{fontenc}
\usepackage{lmodern}
\usepackage{microtype}
\usepackage{enumitem}
\usepackage{booktabs}
\usepackage{array}
\usepackage{graphicx}

\newtheorem{definition}{Definition}[section]
\newtheorem{proposition}[definition]{Proposition}
\newtheorem{theorem}[definition]{Theorem}
\newtheorem{lemma}[definition]{Lemma}
\theoremstyle{remark}
\newtheorem{remark}[definition]{Remark}
\newtheorem{example}[definition]{Example}

\title{Bone Stacking Combinatorial Number System}
\author{Clayton Seelenmayer}
\date{October 4th, 2026}

\begin{document}

\maketitle

\begin{abstract}
This paper defines a finite number system by enumerating combinatorial structures built from at most twelve abstract bones. A structure is first classified by its total number of bones, then by its structural partition, and finally by the orientations of its one-bone, two-bone, and three-bone piles. The resulting system contains exactly \(2{,}911\) distinct values. We derive the counting formula, give ranking and unranking procedures, and provide both concrete and generalized examples.
\end{abstract}

\begin{figure}[h]
	\centering
	\includegraphics[width=0.5\textwidth]{bones.png}
	\caption{Bone number C1, representing decimal 9.}
	\label{fig:myimage}
\end{figure}

\newpage

\section{Introduction}

The system was inspired by the observation that many small mammals have twelve pairs of ribs, as well as a curiousity surrounding the invention of roman numerals. The mathematical construction, however, treats the bones as abstract combinatorial objects. Up to twelve bones are used to represent a number, and are interpreted in a base-2911 number system to allow for individual digits to be read regardless of the positioning of the bones.

\bigskip 
The central idea is to enumerate every distinct structure in a fixed order over three levels:

\begin{enumerate}
	\item Structures are grouped by total number of bones: all one-bone structures,	then all two-bone structures, and so on through twelve bones. This is to minimize the number of bones needed for small numbers, while expanding the total number of representable numbers.
	\item Within a fixed bone-count, structures are grouped by their structural partition. For example, three bones can be partitioned as
	\[
		1+1+1,\qquad 2+1,\qquad 3.
	\]
	Because the piles are unordered, \(1+2\) is not a separate partition from
	\(2+1\).
	\item Within a fixed partition, the possible orientations are enumerated
	according to a fixed canonical ordering.
\end{enumerate}

This gives a one-to-one correspondence between valid bone structures and the natural numbers.
\[
1,2,\ldots,2911.
\]

\newpage
\section{Definitions}
\bigskip 
\begin{definition}[Pile]
A \emph{pile} is a connected group of one, two, or three bones. Four bones is likely to topple over and is not considered. Bone piles can be represented by the alphabet:
\emph{1,A,B,C,D,E,F,G,H,I}.
\end{definition}

\begin{definition}[Bone orientation]
A one-bone pile has one orientation of bone stackings:
\[
\begin{array}{ccl}
1 &=& \text{No connections.}.
\end{array}
\]

A two-bone pile has three orientations:
\[
\begin{array}{ccl}
A &=& \text{end-to-end},\\
B &=& \text{end-to-middle},\\
C &=& \text{middle-to-middle}.
\end{array}
\]

A three-bone pile has six orientations:
\[
\begin{array}{cclcl}
D &=& \text{AA} &=& \text{end-to-end} + \text{end-to-end},\\
E &=& \text{AB} &=& \text{end-to-end} + \text{end-to-middle},\\
F &=& \text{AC} &=& \text{end-to-end} + \text{middle-to-middle},\\
G &=& \text{BB} &=& \text{end-to-middle} + \text{end-to-middle},\\
H &=& \text{BC} &=& \text{end-to-middle} + \text{middle-to-middle},\\
I &=& \text{CC} &=& \text{middle-to-middle} + \text{middle-to-middle}.
\end{array}
\]

The ordering convention is
\[
A<B<C,\qquad D<E<F<G<H<I.
\]
\end{definition}

\begin{definition}[Structural partition]
Suppose a structure consists of \(a\) one-bone piles, \(b\) two-bone piles, and \(c\) three-bone piles, where \(a,b,c\in\mathbb{N}_0\) and \(a+b+c\leq 12\). The \emph{structural partition} of a structure containing \(n\) bones is the triple \((a,b,c)\) satisfying
\[
	a+2b+3c=n.
\]

The piles are treated as unordered within each size class. Thus, for example, \(ABC\) represents the same structure regardless of the order in which the three two-bone piles are listed. The canonical representation is obtained by sorting orientations according to
\[
A<B<C
\]
for two-bone piles and
\[
D<E<F<G<H<I
\]
for three-bone piles.

\bigskip
When different pile sizes occur, three-bone piles are written first, followed by two-bone piles, followed by one-bone piles. Within a partition, smaller letters are written first.\end{definition}

\begin{definition}[Bone-number]
A \emph{bone-number} is the canonical symbolic representation of a valid structure. The decimal value of a bone-number is its position in the global enumeration described above.
\end{definition}

\newpage
\section{The Beginning of the Enumeration}

The first few values are

\[
\begin{array}{rclcrcl}
1&=&1,	&& 9&=&C1,\\
2&=&11,	&& 10&=&D,\\
3&=&A,	&& 11&=&E,\\
4&=&B,	&& 12&=&F,\\
5&=&C,	&& 13&=&G,\\
6&=&111,	&& 14&=&H,\\
7&=&A1,	&& 15&=&I,\\
8&=&B1,	&&
\end{array}
\]

Continuing,

\[
\begin{array}{rclcrcl}
16&=&1111, && 24&=&BC1,\\
17&=&A11,  && 25&=&CC1,\\
18&=&B11,  && 26&=&D1,\\
19&=&C11,  && 27&=&E1,\\
20&=&AA1,  && 28&=&F1,\\
21&=&AB1,  && 29&=&G1,\\
22&=&AC1,  && 30&=&H1,\\
23&=&BB1,  && 31&=&I1.
\end{array}
\]

At the upper end of the system,

\[
\begin{array}{rcl}
2907&=&HHHH,\\
2908&=&HHHI,\\
2909&=&HHII,\\
2910&=&HIII,\\
2911&=&IIII.
\end{array}
\]

\bigskip
Thus the largest value is
\[
\boxed{2911},
\]

\newpage
\section{Counting the Structures}

The counting can be reduced to a simple multiset problem.

\begin{lemma}
For a structural partition containing \(b\) two-bone piles, the number of possible
orientation multisets is
\[
\binom{b+2}{2}.
\]
\end{lemma}

\begin{proof}
There are three possible two-bone orientations, \(A,B,C\). We must choose a
multiset of size \(b\) from these three types. Equivalently, let
\(x_A,x_B,x_C\geq0\) satisfy
\[
x_A+x_B+x_C=b.
\]
By stars and bars, the number of solutions is
\[
\binom{b+3-1}{3-1}
=
\binom{b+2}{2}.
\]
\end{proof}

\begin{lemma}
For a structural partition containing \(c\) three-bone piles, the number of
possible orientation multisets is
\[
\binom{c+5}{5}.
\]
\end{lemma}

\begin{proof}
There are six possible three-bone orientations,
\[
D,E,F,G,H,I.
\]
A multiset of size \(c\) is determined by nonnegative integers
\[
x_D+x_E+x_F+x_G+x_H+x_I=c.
\]
Again by stars and bars, the number of solutions is
\[
\binom{c+6-1}{6-1}
=
\binom{c+5}{5}.
\]
\end{proof}

\begin{proposition}[Structures in a fixed partition]
If a partition contains \("a"\) one-bone piles, \("b"\) two-bone piles, and
\("c"\) three-bone piles, then the number of distinct structures is
\[
\boxed{
\binom{b+2}{2}\binom{c+5}{5}
}.
\]
\end{proposition}

\begin{proof}
The one-bone piles have only one possible orientation, since all one-bone piles are identical. The two-bone piles contribute
\(\binom{b+2}{2}\) orientations, and the three-bone piles contribute
\(\binom{c+5}{5}\) orientations. These choices are independent, so the total is
their product.
\end{proof}

\section{Counting an Entire Bone Block}

Let \(B_n\) denote the number of structures containing exactly \(n\) bones.

\bigskip
For each pair \((b,c)\) satisfying

$$
2b+3c\leq n,
$$

where \(b\) is the number of 2-bone structures and \(c\) is the number of 3-bone structures, the remaining number of bones is

$$
a=n-2b-3c.
$$

We first fix the number \(c\) of 3-bone structures and consider every possible value of \(c\), from \(0\) up to \(\lfloor n/3\rfloor\) since each multiple of 3-bone structures uses exactly 3 bones from the total \(n\). For each fixed \(c\), we then consider every possible number \(b\) of 2-bone structures that can be formed from the remaining bones, from \(0\) up to \(\lfloor(n-3c)/2\rfloor\).

\bigskip
For each pair \((b,c)\), there are

$$
\binom{b+2}{2}
$$

ways to distribute the \(b\) 2-bone structures among their three possible types, and

$$
\binom{c+5}{5}
$$

ways to distribute the \(c\) 3-bone structures among their six possible types.

\bigskip
Therefore,

$$
\boxed{
B_n=
\sum_{c=0}^{\lfloor n/3\rfloor}
\sum_{b=0}^{\lfloor (n-3c)/2\rfloor}
\binom{b+2}{2}\binom{c+5}{5}
}.
$$

\bigskip
Denote the cumulative number of structures through \(n\) bones.
\[
C_n=\sum_{k=1}^{n}B_k
\]

\begin{theorem}[Finite enumeration]
The construction defines exactly \(2911\) distinct bone-numbers.
\end{theorem}

\begin{proof}
Every valid structure is uniquely defined by a total bone count \(n\in\{1,\ldots,12\}\),
a structural partition \((a,b,c)\), and a canonical multiset of
orientations within each pile size. The enumeration orders these disjoint
categories without overlap. Hence each structure occurs exactly once.

\bigskip
The number of structures in the \(n\)-bone block is \(B_n\), so the total number
of structures is
\[
\sum_{n=1}^{12}B_n=C_{12}=2911.
\]
\end{proof}

\section{Ranking and Unranking}

The conversion between decimal integers and bone-numbers can be understood as a
pair of inverse operations.

\subsection{Decimal to Bone-Number}

Let \(N\in\{1,\ldots,2911\}\). First find the unique \(n\) satisfying
\[
C_{n-1}<N\leq C_n,
\]
where \(C_0=0\). The number lies in the \(n\)-bone block. Its one-indexed position within that block is
\[
\boxed{
r=N-C_{n-1}.
}
\]

Next enumerate the structural partitions of \(n\) bones. For each partition, add
\[
\binom{b+2}{2}\binom{c+5}{5}
\]
until the partition containing position \(r\) is found. If the cumulative number before the selected partition is \(P\), then its one-indexed position within that partition is
\[
\boxed{
s=r-P.
}
\]

Finally, unrank the orientation multisets within the selected partition. In a partition containing \(b\) two-bone piles and \(c\) three-bone piles, define
\[
M_2(b)=\binom{b+2}{2},
\qquad
M_3(c)=\binom{c+5}{5}.
\]
The partition therefore contains
\[
M_3(c)M_2(b)
\]
structures. Structures are ordered by first enumerating the \(c\)-multiset of three-bone orientations, then the \(b\)-multiset of two-bone orientations, then the trivial \(a\)-multiset one-bone orientation. The one-indexed ranks of these two components are
\[
\boxed{
p=\left\lfloor\frac{s-1}{M_2(b)}\right\rfloor+1
}
\]
and
\[
\boxed{
q=(s-1)\bmod M_2(b)+1.
}
\]
Thus \(p\) identifies the \(p\)-th three-bone orientation multiset, while \(q\) identifies the \(q\)-th two-bone orientation multiset.

\subsection{Decimal to Bone-Number Example: \(2500\)}

Consider the decimal number
\[
2500.
\]

Since
\[
1847<2500\leq2911,
\]
the number lies in the twelve-bone block. Its position inside that block is

\[
2500-1847=653.
\]

The relevant partition table begins

\[
\begin{array}{l|r|r}
\text{Partition} & \text{Number of structures} & \text{Cumulative}\\
\midrule
1^{12} & 1 & 1\\
2+1^{10} & 3 & 4\\
2^2+1^8 & 6 & 10\\
2^3+1^6 & 10 & 20\\
2^4+1^4 & 15 & 35\\
2^5+1^2 & 21 & 56\\
2^6 & 28 & 84\\
3+1^9 & 6 & 90\\
3+2+1^7 & 18 & 108\\
3+2^2+1^5 & 36 & 144\\
3+2^3+1^3 & 60 & 204\\
3+2^4+1 & 90 & 294\\
3^2+1^6 & 21 & 315\\
3^2+2+1^4 & 63 & 378\\
3^2+2^2+1^2 & 126 & 504\\
3^2+2^3 & 210 & 714
\end{array}
\]

Because
\[
504<653\leq714,
\]
the required partition is
\[
3^2+2^3.
\]

Its local rank is
\[
653-504=149.
\]

Each choice of the two three-bone orientations is followed by ten choices of
the three two-bone orientations. Hence

\[
g
=
\left\lfloor\frac{149-1}{10}\right\rfloor+1
=
15
\]
and
\[
q
=
(149-1)\bmod10+1
=
9.
\]

The \(15\)-th three-bone orientation is
\[
FI,
\]
while the \(9\)-th two-bone orientation is
\[
BCC.
\]

Therefore
\[
\boxed{2500=FIBCC}.
\]

\subsection{Bone-Number to Decimal}

The reverse process is a ranking operation. Given a bone-number, its corresponding decimal value is determined by finding its position in the complete enumeration of structures.

\bigskip
For a general bone-number, proceed as follows:
\begin{enumerate}
	\item Count the total number of bones \(n\).
	\item Determine its structural partition \((a,b,c)\), where \(a\), \(b\),
	      and \(c\) denote the numbers of one-, two-, and three-bone piles,
	      respectively.
	\item Add the complete blocks containing fewer than \(n\) bones. If
	      \(C_{n-1}\) denotes the cumulative number of structures through
	      \(n-1\) bones, this contributes \(C_{n-1}\).
	\item Within the \(n\)-bone block, enumerate the structural partitions in
	      the prescribed order and add the sizes of all partitions preceding
	      the selected partition. Let this cumulative number be \(P\).
	\item Determine the one-indexed rank \(s\) of the given orientation
	      multiset within the selected partition.
\end{enumerate}

If \(P\) is the number of structures preceding the selected partition and
\(s\) is the one-indexed rank within that partition, then the corresponding
decimal value is
\[
\boxed{
N=C_{n-1}+P+s.
}
\]

\subsection{Bone-Number to Decimal Example: \("IX"\)}

See Figure 1 for an example illustration.

\bigskip
Let the observed bone placements \(I\) and \(X\) denote the one-, two-bone orientations respectively:
\[
I=1,
\]
\[
X=C.
\]

Since bone placement permutations are considered equivalent, we use the canonical representation
\[
IX=C1.
\]

This has structural partition
\[
2+1,
\]

Which contains three bones. The structures with fewer than three bones are
\[
11,A,B,C.
\]

Thus the cumulative number of structures through two bones is
\[
C_2=5.
\]

Within the three-bone block, the structural partitions preceding \(2+1\)
consist only of
\[
1+1+1.
\]

This partition contains exactly one structure:
\[
111.
\]

Therefore,
\[
P=1.
\]

Within the \(2+1\) partition, the possible two-bone orientations are
\[
A1,B1,C1.
\]

Thus \(C1\) is the third structure in this partition, so
\[
s=3.
\]

Therefore the global rank is
\[
N=C_2+P+s
=5+1+3
=9.
\]

Hence
\[
\boxed{"IX"=C1=9}.
\]

\section{A Flappy Bird Example}

\subsection{Encoding Game Stat and Program}

The system can also be used as an abstract representation for the state of a
simple game such as \emph{Flappy Bird}. For example, one could assign bone-numbers
to quantities such as

\begin{enumerate}
	\item \(X\)-coordinate,
	\item \(Y\)-coordinate,
	\item velocity,
	\item gravity,
	\item flap impulse,
	\item pipe \(X\)-positions,
	\item pipe gap \(Y\)-positions,
	\item score,
	\item game speed,
	\item random-number state.
\end{enumerate}

Each state variable could be encoded using one or more values in
the range \(1,\ldots,2911\).

\bigskip
A deliberately conservative estimate of roughly \(190\) instructions and \(10\) states would require
\[
200\log_2(2911)
\approx2301
\]
bits, or

\[
\frac{2301}{8}\approx288\text{ bytes}.
\]

This estimate represents the amount of memory needed to store game states and program instructions while excluding graphics, input/output, and other machine states. Since small mammals are imagined to provide two racks of ribs, this would correspond to 100 small mammals.

\subsection{Encoding a \(20\times20\) Monochrome Image}

A \(20\times20\) monochrome image contains
\[
20\cdot20=400
\]
binary pixels, and therefore has
\[
2^{400}
\]
possible states.

\bigskip
Since one bone-number has \(2911\) possible values, the minimum number \(m\) of
bone-numbers required satisfies
\[
2911^m\geq2^{400}.
\]

Taking logarithms gives
\[
m\geq\frac{400}{\log_2(2911)}
\approx34.76.
\]

Hence
\[
\boxed{m=35}.
\]

So \(35\) bone-numbers are sufficient to encode every possible \(20\times20\) monochrome image. Since each small mammal is imagined to provide two racks of ribs, this would correspond to 18 small mammals

\newpage
\section{Conclusion}

The proposed number system is a finite combinatorial enumeration. A value is
identified by the hierarchy

\[
\boxed{
\text{total bone count}
\longrightarrow
\text{structural partition}
\longrightarrow
\text{orientation multiset}.
}
\]

\bigskip
The number of structures for a fixed partition follows from multisets and stars
and bars, giving

\[
\binom{b+2}{2}\binom{c+5}{5}.
\]

\bigskip
Summing over all partitions with at most twelve bones produces exactly
\[
\boxed{2911}
\]
distinct values.

\bigskip
The system therefore provides a well-defined bijection between the natural numbers \(1,\ldots,2911\) and a finite family of canonical bone structures. Decimal integers can be \emph{unranked} into bone-numbers by locating their block, partition, and orientation rank; bone-numbers can be \emph{ranked} back into integers by reversing the same sequence.

\section{Take Away}

Given a set of 12 ribs from a small mammal, you could arrange the bones in piles of 1, 2, and 3, dispersed over a small region, such that you represent the numbers from 1 to 2911. This alphabet of 2911 numerical values is made of only 10 primative symbols, limited structurally, with equivalent permutations.

\end{document}
